\documentclass[../../数学分析培优讲义.tex]{subfiles}

\begin{document}

\setcounter{chapter}{12}
\pagestyle{fancy}

\chapter{曲线积分与曲面积分}


\section{曲线积分}


\subsection{第一型曲线积分}

\begin{theorem}[\markstar\ 计算公式]
设 $L:\bm r=\bm r(t)$, $t\in\lbrack\alpha,\beta\rbrack$ 是一条连续可微的参数曲线, 满足 $\bm r'(t)\ne\bm0$ $(t\in(\alpha,\beta))$, 并设 $f$ 在 $L$ 上连续, 则 $f$ 沿 $L$ 的第一型曲线积分存在, 且
\[
    \int\limits_L f(P)\,\mathrm{d}s
    =\int_\alpha^\beta f\bigl(x(t),y(t),z(t)\bigr)
    \sqrt{[x'(t)]^2+[y'(t)]^2+[z'(t)]^2}\,\mathrm{d}t.
\]
\end{theorem}

因此, 如何求曲线的参数方程十分重要. 对平面曲线
\[
    F(x,y)=0,
\]
常用方法如下: 若能解出 $y=f(x)$ 或 $x=g(y)$, 可直接取 $x$ 或 $y$ 为参数; 也可代入极坐标 $x=r\cos\theta$, $y=r\sin\theta$ 后解出 $r=f(\theta)$ 或 $\theta=g(r)$; 对适合的曲线还可令 $y=tx$ 并求交点坐标作为参数方程.

对于空间曲线
\[
    F(x,y,z)=0,\qquad G(x,y,z)=0,
\]
若能解出其中两个变量为第三个变量的函数, 可直接参数化; 也可代入球坐标或柱坐标, 或先消去一个变量得到平面投影曲线, 再恢复第三个坐标.

\begin{theorem}[\markstar\ 对称性]
若 $L$ 可划分为互相对称的两部分 $L_1,L_2$, 且在对称点上 $f(P)$ 大小相等、符号相反, 则 $L_1,L_2$ 上的积分相互抵消; 若大小相等且符号相同, 则 $L$ 上的积分等于 $L_1$ 上积分的 $2$ 倍.
\end{theorem}

\begin{remark}
区别于重积分, 曲线积分还可利用轮换对称性. 若曲线关于 $x,y,z$ 具有轮换对称性, 则对应不同变量的同形式曲线积分相等.
\end{remark}

\begin{example}
设
\[
    L:\ x^2+y^2+z^2=a^2,\qquad x+y+z=0,
\]
计算积分:
\begin{enumerate}[label=(\arabic*)]
    \item $\displaystyle\int\limits_L x^2\,\mathrm{d}s$;
    \item $\displaystyle\int\limits_L(xy+yz+zx)\,\mathrm{d}s$.
\end{enumerate}
\end{example}
\exampleblank

\begin{example}
计算下列第一型曲线积分:
\begin{enumerate}[label=(\arabic*)]
    \item $\displaystyle I=\int\limits_L\sqrt{x^2+y^2}\,\mathrm{d}s$, 其中 $L:x^2+y^2=ax$ $(a>0)$;
    \item $\displaystyle I=\int\limits_L\left(x^{\frac{4}{3}}+y^{\frac{4}{3}}\right)\,\mathrm{d}s$, 其中 $L:x^{\frac{2}{3}}+y^{\frac{2}{3}}=a^{\frac{2}{3}}$ $(a>0)$.
\end{enumerate}
\end{example}
\exampleblank

\begin{example}
计算
\[
    \int\limits_L(x+y)\,\mathrm{d}s,
\]
其中 $L$ 为顶点 $(0,0),(1,0),(0,1)$ 的三角形边界.
\end{example}
\exampleblank



\subsection{第二型曲线积分}

\begin{theorem}[\markstar]
设 $\bm r(t)=(x(t),y(t),z(t))$, $\alpha\leq t\leq\beta$, 是逐段光滑有向曲线, $P,Q,R$ 在曲线上连续, 则
\[
\begin{aligned}
    \int\limits_L P\,\mathrm{d}x+Q\,\mathrm{d}y+R\,\mathrm{d}z
    ={}&\int_\alpha^\beta\bigl\lbrack P(x(t),y(t),z(t))x'(t)\\
    &\quad+Q(x(t),y(t),z(t))y'(t)+R(x(t),y(t),z(t))z'(t)\bigr\rbrack\,\mathrm{d}t.
\end{aligned}
\]
\end{theorem}

\begin{theorem}[\markstar\ 对称性]
对第二型曲线积分, 除被积函数的大小与符号外, 还需考虑投影元素的符号. 例如对 $\displaystyle\int\limits_L f(P)\,\mathrm{d}x$, 若对称点上 $f(P)\,\mathrm{d}x$ 取相反符号, 则积分为零; 若取相同符号, 则积分等于一半曲线积分的 $2$ 倍.
\end{theorem}

\begin{theorem}[\markstar\ 第一、二型曲线积分的关系]
设 $\bm t,\bm n$ 分别为平面曲线的单位切向量与单位法向量, 则
\[
    \cos(\bm n,x)\,\mathrm{d}s=\mathrm{d}y,
    \qquad
    \cos(\bm n,y)\,\mathrm{d}s=-\mathrm{d}x,
\]
从而
\[
\begin{aligned}
    \int\limits_L P\,\mathrm{d}x+Q\,\mathrm{d}y
    &=\int\limits_L\bigl\lbrack P\cos(\bm t,x)+Q\cos(\bm t,y)\bigr\rbrack\,\mathrm{d}s\\
    &=\int\limits_L\bigl\lbrack -P\cos(\bm n,y)+Q\cos(\bm n,x)\bigr\rbrack\,\mathrm{d}s.
\end{aligned}
\]
\end{theorem}

空间第二型曲线积分常用三种方法: 参数方程化为定积分; Stokes 公式化为第二型曲面积分; 当满足路径无关条件 $P_y=Q_x$, $P_z=R_x$, $Q_z=R_y$ 时取特殊路径或求原函数.

\begin{example}
求曲线积分
\[
    \int\limits_C e^{-(x^2+y^2)}\bigl\lbrack \cos(2xy)\,\mathrm{d}x+\sin(2xy)\,\mathrm{d}y\bigr\rbrack,
\]
其中 $C:x^2+y^2=1$, 方向为逆时针.
\end{example}
\exampleblank

\begin{example}
利用对称性证明
\[
    I=\int\limits_L\dfrac{-x\,\mathrm{d}x+y\,\mathrm{d}y}{x^2+y^2}=0,
\]
其中 $L:x^2+y^2=1$, 方向为逆时针.
\end{example}
\exampleblank

\begin{example}
计算积分
\[
    \int\limits_{ABC}\dfrac{\mathrm{d}x+\mathrm{d}y}{|x|+|y|},
\]
其中 $ABC$ 为三点 $A(1,0),B(0,1),C(-1,0)$ 连成的折线.
\end{example}
\exampleblank

\begin{example}
计算
\[
    \int\limits_L y\,\mathrm{d}x+z\,\mathrm{d}y+x\,\mathrm{d}z,
\]
其中
\[
    L:\ \dfrac{x^2}{a^2}+\dfrac{y^2}{b^2}+\dfrac{z^2}{c^2}=1,
    \qquad \dfrac{x}{a}+\dfrac{z}{c}=1,
    \qquad x,y,z\geq0,
\]
积分方向从 $(a,0,0)$ 到 $(0,0,c)$.
\end{example}
\exampleblank

\begin{example}
计算第一型曲线积分
\[
    \int\limits_C\left(
    \dfrac{xy}{ab}+\dfrac{\sqrt{2}yz}{b\sqrt{a^2+b^2}}+
    \dfrac{\sqrt{2}zx}{a\sqrt{a^2+b^2}}
    \right)\,\mathrm{d}s,
\]
其中 $C$ 为
\[
    \dfrac{x^2}{a^2}+\dfrac{y^2}{b^2}+\dfrac{2z^2}{a^2+b^2}=1,
    \qquad \dfrac{x}{a}+\dfrac{y}{b}=1,
    \qquad x,y,z>0
\]
的交线.
\end{example}
\exampleblank

\begin{example}
证明
\[
    \lim\limits_{R\to+\infty}
    \int\limits_{x^2+y^2=R^2}
    \dfrac{y\,\mathrm{d}x-x\,\mathrm{d}y}{|x^2+xy+y^2|^2}=0.
\]
\end{example}
\exampleblank

\begin{example}
计算
\[
    \int\limits_{\widehat{OA}}(x^2+y^2)\,\mathrm{d}x+4xy\,\mathrm{d}y,
\]
其中:
\begin{enumerate}[label=(\arabic*)]
    \item $\widehat{OA}$ 为上半圆周 $x^2+y^2=ax$ $(y\geq0)$;
    \item $\widehat{OA}$ 为 $x$ 轴上线段 $OA$.
\end{enumerate}
\end{example}
\exampleblank

\begin{example}
计算
\[
    I=\int\limits_L(y-z)\,\mathrm{d}x+(z-x)\,\mathrm{d}y+(x-y)\,\mathrm{d}z,
\]
其中 $L:x^2+y^2+z^2=a^2$, $x+y+z=0$, 从 $z$ 轴正向看 $L$ 取逆时针方向.
\end{example}
\exampleblank

\begin{example}
计算
\[
    \int\limits_{L^+}x\,\mathrm{d}y-y\,\mathrm{d}x,
\]
其中
\[
    L^+:x^{2n+1}+y^{2n+1}=ax^ny^n,\qquad x\geq0,\ y\geq0,
\]
取逆时针方向.
\end{example}
\exampleblank



\subsection{Green 公式}

\begin{theorem}[\markstar\ Green 公式]
\leavevmode\par\noindent
设 $\Omega\subset\mathbb{R}^2$ 是由有限条分段连续可微曲线围成的闭区域, $P(x,y),Q(x,y)$ 在 $\Omega$ 上连续并具有所需偏导数, 则
\[
    \int\limits_{\partial\Omega}P\,\mathrm{d}x+Q\,\mathrm{d}y
    =\iint\limits_\Omega\left(\dfrac{\partial Q}{\partial x}-\dfrac{\partial P}{\partial y}\right)\,\mathrm{d}x\,\mathrm{d}y.
\]
$\partial\Omega$ 取正向, 即沿边界前进时区域位于左侧.
\end{theorem}

\begin{corollary}[\markstar]
记 $S_\Omega$ 为 $\Omega$ 的面积, 则
\[
    S_\Omega
    =\int\limits_{\partial\Omega}x\,\mathrm{d}y
    =-\int\limits_{\partial\Omega}y\,\mathrm{d}x
    =\dfrac{1}{2}\int\limits_{\partial\Omega}(x\,\mathrm{d}y-y\,\mathrm{d}x).
\]
\end{corollary}

\begin{example}
计算下列积分:
\begin{enumerate}[label=(\arabic*)]
    \item $\displaystyle\int\limits_\Gamma xy^2\,\mathrm{d}y-x^2y\,\mathrm{d}x$, $\Gamma:x^2+y^2=a^2$, 按逆时针方向;
    \item $\displaystyle\int\limits_\Gamma(x+y)\,\mathrm{d}x-(x-y)\,\mathrm{d}y$, $\Gamma:\dfrac{x^2}{a^2}+\dfrac{y^2}{b^2}=1$, 按逆时针方向.
\end{enumerate}
\end{example}
\exampleblank

\begin{example}
计算:
\begin{enumerate}[label=(\arabic*)]
    \item $\displaystyle\int\limits_\Gamma\dfrac{x\,\mathrm{d}y-y\,\mathrm{d}x}{4x^2+y^2}$, $\Gamma:(x-1)^2+y^2=R^2$ $(R>1)$, 按逆时针方向;
    \item $\displaystyle\int\limits_L\dfrac{x\,\mathrm{d}y-y\,\mathrm{d}x}{2(Ax^2+2Bxy+Cy^2)}$, $A,C>0$, $AC-B^2>0$, $L:x^2+y^2=1$, 按逆时针方向;
    \item $\displaystyle\int\limits_L\dfrac{x\,\mathrm{d}y-y\,\mathrm{d}x}{(ax+by)^2+(cx+dy)^2}$, $ad-bc\ne0$, $L:(ax+by)^2+(cx+dy)^2=1$, 按逆时针方向.
\end{enumerate}
\end{example}
\exampleblank

\begin{example}
计算
\[
    \int\limits_C\dfrac{e^y}{x^2+y^2}
    \bigl\lbrack (x\sin x+y\cos x)\,\mathrm{d}x+(y\sin x-x\cos x)\,\mathrm{d}y\bigr\rbrack,
\]
其中 $C:x^2+y^2=1$, 按逆时针方向.
\end{example}
\exampleblank

\begin{example}
设 $C$ 为对称于坐标轴的光滑曲线, 证明
\[
    \int\limits_C(x^3y+e^y)\,\mathrm{d}x+(xy^3+xe^y-\cos x)\,\mathrm{d}y=0.
\]
\end{example}
\exampleblank

\begin{example}
设 $D$ 为第一象限内由 $y=x$, $y=4x$, $xy=1$, $xy=4$ 所围区域, $F'(x)=f(x)$, 证明
\[
    \int\limits_{\partial D}\dfrac{F(xy)}{y}\,\mathrm{d}y
    =\ln2\int_1^4f(u)\,\mathrm{d}u,
\]
其中 $\partial D$ 取逆时针方向.
\end{example}
\exampleblank

\begin{example}
计算开口弧段上的曲线积分:
\begin{enumerate}[label=(\arabic*)]
    \item $\displaystyle\int\limits_{C^+}(-2xe^{-x^2}\sin y)\,\mathrm{d}x+(e^{-x^2}\cos y+x^4)\,\mathrm{d}y$, 其中 $C^+$ 为从 $(1,0)$ 到 $(-1,0)$ 的上半圆 $y=\sqrt{1-x^2}$;
    \item $\displaystyle\int\limits_{\widehat{OA}}(y^2-\cos y)\,\mathrm{d}x+x\sin y\,\mathrm{d}y$, 其中 $\widehat{OA}$ 为从 $O(0,0)$ 到 $A(\pi,0)$ 的弧 $y=\sin x$.
\end{enumerate}
\end{example}
\exampleblank

\begin{example}
设 $L$ 为平面上逐段光滑闭曲线, $D$ 为 $L$ 所围区域, $u=u(x,y)$ 在 $D$ 内直到边界有二阶连续偏导数, 试证
\[
    \int\limits_L\dfrac{\partial u}{\partial\bm n}\,\mathrm{d}s
    =\iint\limits_D\left(\dfrac{\partial^2u}{\partial x^2}+\dfrac{\partial^2u}{\partial y^2}\right)\,\mathrm{d}x\,\mathrm{d}y,
\]
其中 $\bm n$ 为 $L$ 的外法向量.
\end{example}
\exampleblank



\subsection{理论证明与应用}

第一型曲线积分与重积分具有许多类似性质, 但第二型曲线积分的性质有明显差别.

\begin{example}
举例说明对第二型曲线积分, ``积分中值定理'' 不再成立.
\end{example}
\exampleblank

\begin{example}
设 $P,Q,R$ 在 $L$ 上连续, $L$ 为光滑弧段, 弧长为 $l$, 证明
\[
    \left|\int\limits_LP\,\mathrm{d}x+Q\,\mathrm{d}y+R\,\mathrm{d}z\right|\leq Ml,
\]
其中
\[
    M=\max_{(x,y,z)\in L}\sqrt{P^2+Q^2+R^2}.
\]
\end{example}
\exampleblank

\begin{example}
证明
\[
    \int\limits_L\cos(\bm l,\bm n)\,\mathrm{d}s=0,
\]
其中 $L$ 为逐段光滑封闭曲线, $\bm l$ 为任意给定方向, $\bm n$ 是 $L$ 的外法线方向.
\end{example}
\exampleblank

\begin{example}
计算 Gauss 积分
\[
    \int\limits_L\dfrac{\cos(\bm r,\bm n)}{r}\,\mathrm{d}s,
\]
其中 $r=\sqrt{(\xi-x)^2+(\eta-y)^2}$, $\bm r$ 为从点 $A(x,y)$ 指向封闭光滑曲线 $L$ 上动点 $M(\xi,\eta)$ 的向量.
\end{example}
\exampleblank

\begin{example}[\markstar]
设 $\bm F=P\bm i+Q\bm j$ 在开区域 $D$ 内处处连续可微, 在 $D$ 内任一圆周 $C$ 上有
\[
    \int\limits_C\bm F\cdot\bm n\,\mathrm{d}s=0,
\]
其中 $\bm n$ 为圆周的单位外法向量, 试证
\[
    \dfrac{\partial P}{\partial x}+\dfrac{\partial Q}{\partial y}=0
\]
在 $D$ 内恒成立.
\end{example}
\exampleblank

\begin{example}[\markstar]
设 $P,Q$ 在全平面上有连续偏导数, 且以任意点 $(x_0,y_0)$ 为中心、任意 $r>0$ 为半径的上半圆 $C$ 上恒有
\[
    \int\limits_C P(x,y)\,\mathrm{d}x+Q(x,y)\,\mathrm{d}y=0,
\]
证明
\[
    P(x,y)\equiv0,
    \qquad
    \dfrac{\partial Q}{\partial y}\equiv0.
\]
\end{example}
\exampleblank

\begin{example}
计算
\[
    \int\limits_{x^2+y^2=R^2}\ln\sqrt{(x-a)^2+(y-b)^2}\,\mathrm{d}s.
\]
\end{example}
\exampleblank

\begin{example}
设 $D:x^2+y^2\leq a^2$ $(a>0)$, $f(x,y)$ 在 $D$ 上有连续偏导数, 且 $f=0$ 在 $\partial D$ 上成立, 证明
\[
    \left|\iint\limits_Df(x,y)\,\mathrm{d}x\,\mathrm{d}y\right|
    \leq\dfrac{\pi a^3}{3}\max_{(x,y)\in D}
    \sqrt{\left(\dfrac{\partial f}{\partial x}\right)^2+
    \left(\dfrac{\partial f}{\partial y}\right)^2}.
\]
\end{example}
\exampleblank



\methodsection{(A) 调和函数 *}

\begin{theorem}[\markstar]
设闭区域 $D$ 由有限条逐段光滑曲线围成, $u,v\in C^{(2)}(D)$, 则
\[
    \iint\limits_D\Delta u\,\mathrm{d}x\,\mathrm{d}y
    =\int\limits_{\partial D}\dfrac{\partial u}{\partial\bm n}\,\mathrm{d}s,
\]
\[
    \iint\limits_Dv\Delta u\,\mathrm{d}x\,\mathrm{d}y
    =-\iint\limits_D\left(
    \dfrac{\partial u}{\partial x}\dfrac{\partial v}{\partial x}
    +\dfrac{\partial u}{\partial y}\dfrac{\partial v}{\partial y}
    \right)\,\mathrm{d}x\,\mathrm{d}y
    +\int\limits_{\partial D}v\dfrac{\partial u}{\partial\bm n}\,\mathrm{d}s,
\]
以及 Green 第二公式
\[
    \iint\limits_D(v\Delta u-u\Delta v)\,\mathrm{d}x\,\mathrm{d}y
    =\int\limits_{\partial D}\left(v\dfrac{\partial u}{\partial\bm n}-u\dfrac{\partial v}{\partial\bm n}\right)\,\mathrm{d}s.
\]
\end{theorem}

\begin{example}
设 $L$ 为平面上一条自身不相交的光滑曲线, 起点 $(1,0)$, 终点 $(0,2)$, 除起终点外全部位于第一象限, 计算
\[
    \int\limits_L\dfrac{\partial\ln r}{\partial\bm n}\,\mathrm{d}s,
\]
其中 $r$ 表示 $L$ 上动点到原点的距离.
\end{example}
\exampleblank

\begin{example}
设 $\Omega$ 为 $xOy$ 平面上具有光滑边界的有界闭区域, $u$ 在 $\Omega$ 内有二阶连续偏导数, 直到边界有一阶连续偏导数, 且 $u|_{\partial\Omega}=0$, $u$ 非常值, 证明
\[
    \iint\limits_\Omega u\left(\dfrac{\partial^2u}{\partial x^2}+\dfrac{\partial^2u}{\partial y^2}\right)\,\mathrm{d}x\,\mathrm{d}y<0.
\]
\end{example}
\exampleblank

\begin{example}
设 $u(x,y)$ 有二阶连续偏导数, 证明: $u$ 为调和函数当且仅当对任意逐段光滑闭曲线 $C$,
\[
    \int\limits_C\dfrac{\partial u}{\partial\bm n}\,\mathrm{d}s=0,
\]
其中 $\bm n$ 为外法线方向.
\end{example}
\exampleblank

\begin{example}
设 $D$ 为区域, $u\in C^{(2)}(D)$, 证明: $u$ 为 $D$ 上调和函数当且仅当对任意 $P_0(x_0,y_0)\in D$,
\[
    u(x_0,y_0)=\dfrac{1}{2\pi}\int_0^{2\pi}u(x_0+r\cos\theta,y_0+r\sin\theta)\,\mathrm{d}\theta,
\]
其中 $0<r<d(P_0,\partial D)$.
\end{example}
\exampleblank

\begin{example}
设 $L$ 是 $\mathbb{R}^2$ 中简单闭曲线, 其所围区域为 $D$, $f$ 在 $D$ 上连续且为调和函数, 证明 $f$ 在 $L$ 上取到最大值和最小值.
\end{example}
\exampleblank

\begin{example}
设闭区域 $D$ 由有限条逐段光滑曲线围成, 调和函数 $u(x,y)$ 在边界 $\partial D$ 上为零, 证明 $u$ 在 $D$ 上恒为零.
\end{example}
\exampleblank

\begin{example}
设 $f(x,y,z)$ 为 $n$ 次齐次函数, $f\in C^{(2)}(\mathbb{R}^3)$, 若对任意球面 $S$ 都有
\[
    \iint\limits_Sf(x,y,z)\,\mathrm{d}S=0,
\]
证明 $\Delta f=0$.
\end{example}
\exampleblank

\begin{example}
设 $\Sigma$ 是 $\mathbb{R}^3$ 第一卦限内任意光滑闭曲面, 所围立体为 $\Omega$, $f(x,y,z)$ 在第一卦限上具有一阶偏导数, 给出
\[
    \iint\limits_\Sigma f(x,y,z)(x\,\mathrm{d}y\,\mathrm{d}z+y\,\mathrm{d}z\,\mathrm{d}x+z\,\mathrm{d}x\,\mathrm{d}y)=0
\]
的充要条件.
\end{example}
\exampleblank



\methodsection{(B) 曲线积分与路径无关性}

\begin{theorem}[\markstar]
若 $P,Q$ 在区域 $D$ 内连续并有连续一阶偏导数, 且 $D$ 为单连通区域, 则下列命题等价:
\begin{enumerate}[label=(\arabic*)]
    \item $\displaystyle\int\limits_LP\,\mathrm{d}x+Q\,\mathrm{d}y$ 只与起点、终点有关;
    \item 对 $D$ 内任意分段光滑闭曲线 $C$, $\displaystyle\int\limits_CP\,\mathrm{d}x+Q\,\mathrm{d}y=0$;
    \item $P\,\mathrm{d}x+Q\,\mathrm{d}y$ 为恰当微分, 即存在 $u$ 使 $u_x=P,u_y=Q$;
    \item 在 $D$ 内处处成立 $\dfrac{\partial Q}{\partial x}=\dfrac{\partial P}{\partial y}$.
\end{enumerate}
\end{theorem}

\begin{remark}
原函数可写成
\[
    u(x,y)=\int\limits_{(x_0,y_0)}^{(x,y)}P\,\mathrm{d}x+Q\,\mathrm{d}y.
\]
求原函数时可选特殊路径, 可逐步求不定积分, 也可利用
\[
    u\,\mathrm{d}v+v\,\mathrm{d}u=\mathrm{d}(uv),
    \qquad
    \dfrac{u\,\mathrm{d}v-v\,\mathrm{d}u}{u^2}=\mathrm{d}\left(\dfrac{v}{u}\right),
\]
\[
    \dfrac{u\,\mathrm{d}v-v\,\mathrm{d}u}{u^2+v^2}=\mathrm{d}\left(\arctan\dfrac{v}{u}\right)
\]
凑全微分. 对多连通区域, 条件 $Q_x=P_y$ 只为必要条件; 各洞对应循环常数, 循环常数全为零时积分才与路径无关.
\end{remark}

\begin{example}
计算
\[
    \int\limits_{\widehat{AB}}(x^2+2xy-y^2)\,\mathrm{d}x+(x^2-2xy-y^2)\,\mathrm{d}y,
\]
其中 $A=(0,0)$, $B=(2,1)$, $\widehat{AB}$ 为任意连接 $A,B$ 的光滑曲线.
\end{example}
\exampleblank

\begin{example}
设 $f(x)$ 在 $(-\infty,+\infty)$ 上有连续导函数, 求
\[
    \int\limits_L
    \dfrac{1+y^2f(x,y)}{y}\,\mathrm{d}x
    +\dfrac{x}{y^2}\bigl\lbrack y^2f(x,y)-1\bigr\rbrack\,\mathrm{d}y,
\]
其中 $L$ 是从 $A(3,\dfrac{2}{3})$ 到 $B(1,2)$ 的直线段.
\end{example}
\exampleblank

\begin{example}
求
\[
    \int\limits_L[ e^x\sin y-3(x+y)]\,\mathrm{d}x+(e^x\cos y-ax)\,\mathrm{d}y,
    \qquad a>0,
\]
其中 $L$ 从 $A=(2a,0)$ 沿 $y=\sqrt{2ax-x^2}$ 到 $B=(0,0)$.
\end{example}
\exampleblank

\begin{example}
计算
\[
    \int\limits_Lx\ln(x^2+y^2-1)\,\mathrm{d}x+y\ln(x^2+y^2-1)\,\mathrm{d}y,
\]
其中 $L$ 是从 $A(2,0)$ 到 $B(0,2)$ 的逐段光滑曲线.
\end{example}
\exampleblank

\begin{example}
计算
\[
    I=\int\limits_{L^+}\dfrac{(x+y)\,\mathrm{d}x-(x-y)\,\mathrm{d}y}{x^2+y^2},
\]
其中 $L^+$ 是从 $A(-1,0)$ 到 $B(1,0)$ 的一条不通过原点的光滑曲线, 满足 $y=f(x)$, $-1\leq x\leq1$.
\end{example}
\exampleblank

\begin{remark}
本题可思考是否能够直接通过求原函数计算积分 $I$.
\end{remark}

\begin{example}
计算下列曲线积分:
\begin{enumerate}[label=(\arabic*)]
    \item $\displaystyle\int\limits_{\widehat{AB}}x\,\mathrm{d}x+y\,\mathrm{d}y$, $A=(0,1)$, $B=(3,4)$;
    \item $\displaystyle\int\limits_{\widehat{AB}}\dfrac{y\,\mathrm{d}x-x\,\mathrm{d}y}{x^2}$, $A=(2,1)$, $B=(1,2)$, $\widehat{AB}$ 不穿过 $y$ 轴;
    \item $\displaystyle\int\limits_{\widehat{AB}}\dfrac{x\,\mathrm{d}x+y\,\mathrm{d}y}{\sqrt{x^2+y^2}}$, $A=(1,0)$, $B=(6,8)$, $\widehat{AB}$ 不过原点.
\end{enumerate}
\end{example}
\exampleblank

\begin{example}
求下列表达式的原函数:
\begin{enumerate}[label=(\arabic*)]
    \item $\displaystyle\dfrac{y\,\mathrm{d}x-x\,\mathrm{d}y}{3x^2-2xy+3y^2}$, $y>0$;
    \item $\displaystyle\dfrac{(x^2+2xy+5y^2)\,\mathrm{d}x+(x^2-2xy+y^2)\,\mathrm{d}y}{(x+y)^2}$.
\end{enumerate}
\end{example}
\exampleblank

\begin{example}
求 $P,Q\in C^{(1)}(\mathbb{R}^2)$ 的关系, 使得对任意光滑闭曲线 $L\subset\mathbb{R}^2$,
\[
    I(\alpha,\beta)=\int\limits_LP(x+\alpha,y+\beta)\,\mathrm{d}x+Q(x+\alpha,y+\beta)\,\mathrm{d}y
\]
与常数 $\alpha,\beta$ 无关.
\end{example}
\exampleblank


\newpage

\section{曲面积分}

本节总体脉络与曲线积分类似, 差别主要体现在维数与计算复杂度.



\subsection{第一型曲面积分}

\begin{theorem}[\markstar]
\begin{enumerate}[label=(\arabic*)]
    \item 设 $S:z=z(x,y)$, $\Delta$ 为 $S$ 在 $xOy$ 平面的投影, 则
    \[
        \iint\limits_S f(x,y,z)\,\mathrm{d}S
        =\iint\limits_\Delta f(x,y,z(x,y))\sqrt{1+z_x^2+z_y^2}\,\mathrm{d}x\,\mathrm{d}y.
    \]
    \item 设
    \[
        S:\begin{cases}
        x=x(u,v),\\
        y=y(u,v),\\
        z=z(u,v),
        \end{cases}
    \]
    记
    \[
        A=\dfrac{\partial(y,z)}{\partial(u,v)},\quad
        B=\dfrac{\partial(z,x)}{\partial(u,v)},\quad
        C=\dfrac{\partial(x,y)}{\partial(u,v)},
    \]
    则
    \[
        \iint\limits_Sf(x,y,z)\,\mathrm{d}S
        =\iint\limits_\Delta f(x(u,v),y(u,v),z(u,v))\sqrt{A^2+B^2+C^2}\,\mathrm{d}u\,\mathrm{d}v.
    \]
\end{enumerate}
\end{theorem}

\begin{remark}
第一型曲面积分通常直接使用计算公式, 或先化为第二型曲面积分再利用第二型曲面积分的方法.
\end{remark}

\begin{example}
设 $f(z)$ 为奇函数, 求
\[
    I=\iint\limits_Sf(z)\,\mathrm{d}S,\qquad
    J=\iint\limits_Sf^2(z)\,\mathrm{d}S,\qquad
    R=\iint\limits_Syf^2(z)\,\mathrm{d}S,
\]
其中 $S$ 为锥面 $z^2=2xy$ 位于球面 $x^2+y^2+z^2=a^2$ 内的部分.
\end{example}
\exampleblank

\begin{example}
计算锥面
\[
    z=\sqrt{2xy},\qquad x\geq0,\ y\geq0,
\]
位于球面 $x^2+y^2+z^2=a^2$ 内部分的面积.
\end{example}
\exampleblank

\begin{example}
计算
\[
    \iint\limits_S|z|\,\mathrm{d}S,
\]
其中 $S$ 为柱体 $x^2+y^2\leq ax$ 被球体 $x^2+y^2+z^2\leq a^2$ 截取部分的表面.
\end{example}
\exampleblank

\begin{example}
计算
\[
    \iint\limits_Sz\,\mathrm{d}S,
\]
其中 $S$ 是曲面 $x^2+z^2=2az$ $(a>0)$ 被曲面 $z=\sqrt{x^2+y^2}$ 所截取的有限部分.
\end{example}
\exampleblank

\begin{example}
计算曲面积分
\[
    F(t)=\iint\limits_{x+y+z=t}f(x,y,z)\,\mathrm{d}S,
\]
其中
\[
    f(x,y,z)=
    \begin{cases}
        1-(x^2+y^2+z^2), & x^2+y^2+z^2\leq1,\\
        0, & x^2+y^2+z^2>1.
    \end{cases}
\]
\end{example}
\exampleblank

\begin{example}
计算
\[
    \iint\limits_\Sigma\dfrac{\mathrm{d}S}{\sqrt{x^2+y^2+(z+a)^2}},
\]
其中 $\Sigma$ 为以原点为中心、$a$ 为半径的上半球面.
\end{example}
\exampleblank

\begin{example}
已知 $S=\{(x,y,z)\mid x^2+y^2+z^2=1\}$, 试用 $\Gamma$ 函数表示
\[
    \iint\limits_S(x^2+y^2)^\beta\,\mathrm{d}S,
    \qquad \beta>-1.
\]
\end{example}
\exampleblank

\begin{example}
设 $f$ 连续, 证明 Poisson 公式
\[
    \int_0^{2\pi}\int_0^\pi
    f(a\sin\varphi\cos\theta+b\sin\varphi\sin\theta+c\cos\varphi)
    \sin\varphi\,\mathrm{d}\varphi\,\mathrm{d}\theta
    =2\pi\int_{-1}^1f(kz)\,\mathrm{d}z,
\]
其中 $k=\sqrt{a^2+b^2+c^2}$.
\end{example}
\exampleblank

\begin{example}
设 $f\in C^{(1)}\lbrack-1,1\rbrack$, $f(-1)=f(1)=0$,
\[
    M=\max_{-1\leq x\leq1}|f'(x)|,
\]
$S$ 为单位球面, 常数 $m,n,p$ 满足 $m^2+n^2+p^2=1$, 证明
\[
    \left|\iint\limits_Sf(mx+ny+pz)\,\mathrm{d}S\right|\leq2\pi M.
\]
\end{example}
\exampleblank

\begin{example}
设 $\Sigma:x^2+y^2+z^2=1$, $A$ 为 $\Sigma$ 内一点, $A$ 到原点的距离为 $q$ $(0<q<1)$, $\rho$ 表示 $A$ 到动点 $P\in\Sigma$ 的距离, 求
\[
    I=\iint\limits_\Sigma\dfrac{\mathrm{d}S}{\rho^2}.
\]
\end{example}
\exampleblank

\begin{example}
设 $S$ 为椭球面, $\rho$ 表示从椭球中心到椭球表面元素 $\mathrm{d}S$ 所在切平面的距离, 计算
\[
    \text{(1) }\iint\limits_S\rho\,\mathrm{d}S;
    \qquad
    \text{(2) }\iint\limits_S\dfrac{1}{\rho}\,\mathrm{d}S.
\]
\end{example}
\exampleblank



\subsection{第二型曲面积分}

求第二型曲面积分常用方法为: Stokes 公式、Gauss 公式以及参数方程. 计算时必须注意投影面积元素的符号.

\begin{theorem}[\markstar]
\[
    \iint\limits_{S^+}P\,\mathrm{d}y\,\mathrm{d}z+Q\,\mathrm{d}z\,\mathrm{d}x+R\,\mathrm{d}x\,\mathrm{d}y
    =\iint\limits_S(P\cos\alpha+Q\cos\beta+R\cos\gamma)\,\mathrm{d}S,
\]
其中单位法向量 $\bm n=(\cos\alpha,\cos\beta,\cos\gamma)$.
\end{theorem}

\begin{example}
设 $f$ 为奇函数, $S^+$ 为 $|x|+|y|+|z|=1$ 的外侧, 利用对称性求出或化简
\[
\begin{aligned}
    I&=\iint\limits_{S^+}(\mathrm{d}x\,\mathrm{d}y+\mathrm{d}y\,\mathrm{d}z+\mathrm{d}z\,\mathrm{d}x),\\
    J&=\iint\limits_{S^+}f^2(z)\,\mathrm{d}x\,\mathrm{d}y,\\
    K&=\iint\limits_{S^+}xf^2(z)\,\mathrm{d}x\,\mathrm{d}y,\\
    L&=\iint\limits_{S^+}f(z)\,\mathrm{d}x\,\mathrm{d}y,\\
    M&=\iint\limits_{S^+}(x+2y+3z)f(x+y+z)\,\mathrm{d}x\,\mathrm{d}y.
\end{aligned}
\]
\end{example}
\exampleblank

\begin{example}
设
\[
    S^+:z-c=\sqrt{R^2-(x-a)^2-(y-b)^2}
\]
取上侧, 计算
\[
    \iint\limits_{S^+}x^2\,\mathrm{d}y\,\mathrm{d}z+y^2\,\mathrm{d}z\,\mathrm{d}x+(x-a)y z\,\mathrm{d}x\,\mathrm{d}y.
\]
\end{example}
\exampleblank

\begin{example}
计算
\[
    \dfrac{1}{3}\iint\limits_S(x\,\mathrm{d}y\,\mathrm{d}z+y\,\mathrm{d}z\,\mathrm{d}x+z\,\mathrm{d}x\,\mathrm{d}y),
\]
其中 $S$ 是球面 $x^2+y^2+z^2=a^2$ 的外侧.
\end{example}
\exampleblank

\begin{example}
设 $\Delta$ 是以 $(1,0,0),(0,1,0),(0,0,1)$ 为顶点的三角形的上侧, 计算
\[
    I=\iint\limits_\Delta(x\,\mathrm{d}y\,\mathrm{d}z+y\,\mathrm{d}z\,\mathrm{d}x+z\,\mathrm{d}x\,\mathrm{d}y),
\]
\[
    J=\iint\limits_\Delta(x^2\,\mathrm{d}y\,\mathrm{d}z+y^2\,\mathrm{d}z\,\mathrm{d}x+z^2\,\mathrm{d}x\,\mathrm{d}y).
\]
\end{example}
\exampleblank

\begin{example}
求
\[
    \iint\limits_Sz^2\cos\gamma\,\mathrm{d}S,
\]
其中 $S$ 为上半球面 $x^2+y^2+z^2=1$, $z\geq0$, $\gamma$ 为球面外法线与 $z$ 轴的夹角.
\end{example}
\exampleblank



\subsection{Gauss 公式和 Stokes 公式}

\begin{example}
计算下列积分:
\begin{enumerate}[label=(\arabic*)]
    \item $\displaystyle\iint\limits_\Sigma(x-y)\,\mathrm{d}y\,\mathrm{d}z+(y-z)\,\mathrm{d}z\,\mathrm{d}x+(z-x)\,\mathrm{d}x\,\mathrm{d}y$, 其中 $\Sigma:z=x^2+y^2$, $z\leq1$, 方向朝下;
    \item $\displaystyle\iint\limits_{S^+}(x+y-z)\,\mathrm{d}y\,\mathrm{d}z+[2y+\sin(z+x)]\,\mathrm{d}z\,\mathrm{d}x+(3z+e^{x+y})\,\mathrm{d}x\,\mathrm{d}y$, 其中 $S^+$ 为 $|x-y+z|+|y-z+x|+|z-x+y|=1$ 的外表面.
\end{enumerate}
\end{example}
\exampleblank

\begin{example}
计算下列积分:
\begin{enumerate}[label=(\arabic*)]
    \item $\displaystyle\iint\limits_Sx^3\,\mathrm{d}y\,\mathrm{d}z+y^3\,\mathrm{d}z\,\mathrm{d}x+z^3\,\mathrm{d}x\,\mathrm{d}y$, 其中 $S$ 是上半椭球面 $\dfrac{x^2}{a^2}+\dfrac{y^2}{b^2}+\dfrac{z^2}{c^2}=1$, $z\geq0$, 方向朝上;
    \item $\displaystyle\iint\limits_\Sigma z\,\mathrm{d}x\,\mathrm{d}y+y\,\mathrm{d}z\,\mathrm{d}x+x\,\mathrm{d}y\,\mathrm{d}z$, 其中 $\Sigma$ 为圆柱面 $x^2+y^2=1$ 被 $z=0,z=3$ 截取部分的外侧.
\end{enumerate}
\end{example}
\exampleblank

\begin{example}
若 $S$ 为封闭光滑曲面, $\bm l$ 为任意固定方向, 证明
\[
    \iint\limits_S\cos(\bm n,\bm l)\,\mathrm{d}S=0.
\]
\end{example}
\exampleblank

\begin{example}
记 $r=r(\theta,\varphi)$ 为分片光滑封闭曲面 $S$ 的球坐标方程, 证明 $S$ 所围区域体积满足
\[
    V=\dfrac{1}{3}\iint\limits_Sr\cos\psi\,\mathrm{d}S,
\]
其中 $\psi$ 为曲面外法线方向与径向量的夹角.
\end{example}
\exampleblank

\begin{example}
设 $V$ 是不含原点的有界闭区域, 边界为光滑简单闭曲面 $\Sigma$, 单位外法向量为 $\bm n$, 径向量 $\bm r=(x,y,z)$, $f\in C^1\lbrack0,+\infty)$, 且
\[
    f'(t)+2f(t)-1=0,
\]
计算
\[
    \iint\limits_\Sigma f\left(\sqrt{x^2+y^2+z^2}\right)\cos(\bm r,\bm n)\,\mathrm{d}S.
\]
\end{example}
\exampleblank

\begin{example}
计算
\[
    \iint\limits_S\dfrac{\cos(\bm n,\bm r)}{r^2}\,\mathrm{d}S,
\]
其中 $S$ 是光滑封闭曲面, 原点不在 $S$ 上, $r$ 为动点到原点的距离.
\end{example}
\exampleblank

\begin{example}
计算
\[
    \iint\limits_{S_{\text{外}}}\dfrac{x\,\mathrm{d}y\,\mathrm{d}z+y\,\mathrm{d}z\,\mathrm{d}x+z\,\mathrm{d}x\,\mathrm{d}y}{(x^2+y^2+z^2)^{\frac{3}{2}}},
\]
其中 $S$ 是立方体 $|x|\leq2$, $|y|\leq2$, $|z|\leq2$ 的表面.
\end{example}
\exampleblank

\begin{example}
设 $V\subset\mathbb{R}^3$ 是其内任一封闭曲面均可在 $V$ 内连续收缩至一点的区域, $P,Q,R$ 在 $V$ 上有连续偏导数. 证明: 对 $V$ 内任一不自交光滑封闭曲面 $S$,
\[
    \iint\limits_S\bigl\lbrack P\cos(\bm n,x)+Q\cos(\bm n,y)+R\cos(\bm n,z)\bigr\rbrack\,\mathrm{d}S=0
\]
当且仅当
\[
    \dfrac{\partial P}{\partial x}+\dfrac{\partial Q}{\partial y}+\dfrac{\partial R}{\partial z}=0
\]
在 $V$ 内处处成立.
\end{example}
\exampleblank

\begin{example}
设 $P(x,y,z),Q(x,y,z),R(x,y,z)$ 在 $\mathbb{R}^3$ 中有一阶偏导数. 对任意 $r>0$ 以及任意点
\[
    (x_0,y_0,z_0)\in\mathbb{R}^3,
\]
考虑半球面
\[
    S:z=z_0+\sqrt{r^2-(x-x_0)^2-(y-y_0)^2}
\]
上的积分
\[
    \iint\limits_S
    P\,\mathrm{d}y\,\mathrm{d}z
    +Q\,\mathrm{d}z\,\mathrm{d}x
    +R\,\mathrm{d}x\,\mathrm{d}y=0.
\]
证明
\[
    \dfrac{\partial P}{\partial x}+\dfrac{\partial Q}{\partial y}=0,
    \qquad R=0
\]
在 $\mathbb{R}^3$ 内处处成立.
\end{example}
\exampleblank

\begin{example}
计算
\[
    \int\limits_{L^+}(z-y)\,\mathrm{d}x+(x-z)\,\mathrm{d}y+(y-x)\,\mathrm{d}z,
\]
其中 $L^+$ 从 $A(a,0,0)$ 经 $B(0,a,0)$ 到 $C(0,0,a)$ 再回到 $A$.
\end{example}
\exampleblank

\begin{example}
计算
\[
    \int\limits_{L^+}(y^2+z^2)\,\mathrm{d}x+(x^2+z^2)\,\mathrm{d}y+(x^2+y^2)\,\mathrm{d}z,
\]
其中 $L$ 是曲面 $x^2+y^2+z^2=4x$ 与 $x^2+y^2=2x$ 的交线中 $z\geq0$ 的部分, 积分方向从原点进入第一卦限.
\end{example}
\exampleblank

\begin{example}
计算
\[
    \int\limits_{L^+}y\,\mathrm{d}x+z\,\mathrm{d}y+x\,\mathrm{d}z,
\]
其中 $L^+$ 为圆周 $x^2+y^2+z^2=a^2$, $x+y+z=0$, 从 $z$ 轴正方向看为逆时针.
\end{example}
\exampleblank

\begin{example}
计算
\[
    \int\limits_{L^+}(x^2-yz)\,\mathrm{d}x+(y^2-xz)\,\mathrm{d}y+(z^2-xy)\,\mathrm{d}z,
\]
其中 $L^+$ 是从 $A(1,0,0)$ 到 $B(1,0,2)$ 的光滑曲线.
\end{example}
\exampleblank

\begin{example}
设 $L$ 为空间某封闭光滑曲线, $P,Q,R$ 为空间中具有一阶连续偏导数的函数, 证明
\[
\left|\int\limits_{L^+}P\,\mathrm{d}x+Q\,\mathrm{d}y+R\,\mathrm{d}z\right|
\leq
\max_{(x,y,z)\in S}
\sqrt{
\left(\dfrac{\partial Q}{\partial x}-\dfrac{\partial P}{\partial y}\right)^2+
\left(\dfrac{\partial R}{\partial y}-\dfrac{\partial Q}{\partial z}\right)^2+
\left(\dfrac{\partial P}{\partial z}-\dfrac{\partial R}{\partial x}\right)^2}
\cdot S,
\]
其中 $S$ 表示以 $L$ 为边界的某曲面, 同时也用 $S$ 表示其面积.
\end{example}
\exampleblank


\newpage

\section{场论初步 *}

记
\[
    \nabla=\left(\dfrac{\partial}{\partial x},\dfrac{\partial}{\partial y},\dfrac{\partial}{\partial z}\right)
\]
为 Nabla 算子. 设 $u(x,y,z)$ 为光滑数量场, $\bm A=P\bm i+Q\bm j+R\bm k$ 为光滑向量场.

\begin{definition}
梯度、散度和旋度分别定义为
\[
    \operatorname{grad}u=\nabla u
    =\left(\dfrac{\partial u}{\partial x},\dfrac{\partial u}{\partial y},\dfrac{\partial u}{\partial z}\right),
\]
\[
    \operatorname{div}\bm A=\nabla\cdot\bm A
    =\dfrac{\partial P}{\partial x}+\dfrac{\partial Q}{\partial y}+\dfrac{\partial R}{\partial z},
\]
\[
    \operatorname{rot}\bm A=\nabla\times\bm A
    =\left(
    \dfrac{\partial R}{\partial y}-\dfrac{\partial Q}{\partial z},
    \dfrac{\partial P}{\partial z}-\dfrac{\partial R}{\partial x},
    \dfrac{\partial Q}{\partial x}-\dfrac{\partial P}{\partial y}
    \right).
\]
\end{definition}

\begin{remark}
若 $\mathrm{d}\bm S=\bm n\,\mathrm{d}S$, $\mathrm{d}\bm r=(\mathrm{d}x,\mathrm{d}y,\mathrm{d}z)$, 则 Gauss 公式和 Stokes 公式可分别写为
\[
    \iint\limits_S\bm A\cdot\mathrm{d}\bm S
    =\iiint\limits_V\operatorname{div}\bm A\,\mathrm{d}V,
\]
\[
    \int\limits_L\bm A\cdot\mathrm{d}\bm r
    =\iint\limits_S(\operatorname{rot}\bm A)\cdot\mathrm{d}\bm S.
\]
\end{remark}

\begin{example}
证明 $\nabla$ 的下列性质:
\begin{enumerate}[label=(\arabic*)]
    \item $\nabla(cf)=c\nabla f$, $c$ 为常数;
    \item $\nabla(f\pm g)=\nabla f\pm\nabla g$;
    \item $\nabla(fg)=f\nabla g+g\nabla f$;
    \item 若 $\varphi$ 为单变量函数, 则 $\nabla(\varphi\circ f)=(\varphi'\circ f)\nabla f$.
\end{enumerate}
\end{example}
\exampleblank

\begin{example}
证明 $\operatorname{div}$ 的下列性质:
\begin{enumerate}[label=(\arabic*)]
    \item $\operatorname{div}(c\bm A)=c\operatorname{div}\bm A$, $c$ 为常数;
    \item $\operatorname{div}(\bm A\pm\bm B)=\operatorname{div}\bm A\pm\operatorname{div}\bm B$;
    \item $\operatorname{div}(u\bm A)=u\operatorname{div}\bm A+\bm A\cdot\nabla u$.
\end{enumerate}
\end{example}
\exampleblank

\begin{example}
证明 $\operatorname{rot}$ 的下列性质:
\begin{enumerate}[label=(\arabic*)]
    \item $\operatorname{rot}(c\bm A)=c\operatorname{rot}\bm A$, $c$ 为常数;
    \item $\operatorname{rot}(\bm A+\bm B)=\operatorname{rot}\bm A+\operatorname{rot}\bm B$;
    \item $\operatorname{rot}(u\bm A)=u\operatorname{rot}\bm A+\nabla u\times\bm A$;
    \item $\nabla\cdot(\bm F_1\times\bm F_2)=(\nabla\times\bm F_1)\cdot\bm F_2-(\nabla\times\bm F_2)\cdot\bm F_1$.
\end{enumerate}
\end{example}
\exampleblank

\begin{definition}
设向量场 $\bm F=(P,Q,R)$ 定义在区域 $D\subseteq\mathbb{R}^3$ 上. 若存在 $D$ 上数量场 $\varphi$, 满足 $\nabla\varphi=\bm F$, 则称 $\bm F$ 为有势场, $\varphi$ 称为 $\bm F$ 的势函数.
\end{definition}

\begin{definition}
设 $\bm F$ 是定义在区域 $D\subseteq\mathbb{R}^3$ 上的向量场. 若对 $D$ 中任意封闭曲线 $\Gamma$ 都有
\[
    \int\limits_\Gamma\bm F\cdot\mathrm{d}\bm r=0,
\]
则称 $\bm F$ 为 $D$ 上的保守场.
\end{definition}

\begin{example}
若 $\operatorname{rot}\bm F=\nabla\times\bm F=0$ 在 $D$ 上处处成立, 则称 $\bm F$ 为 $D$ 上的无旋场.
\end{example}
\exampleblank

\begin{theorem}[\markstar]
设 $D$ 是 $\mathbb{R}^3$ 中的曲面单连通区域, $\bm F\in C^2(D)$, 则下列命题等价:
\begin{enumerate}[label=(\arabic*)]
    \item $\bm F$ 为有势场;
    \item $\bm F$ 为无旋场;
    \item $\bm F$ 为保守场.
\end{enumerate}
\end{theorem}

\begin{example}
证明上述定理.
\end{example}
\exampleblank

\end{document}
